\documentclass[a4paper, noanswers, addpoints, 12pt]{exam}
\usepackage[T1]{fontenc}
\usepackage[utf8]{inputenc}
\usepackage{lmodern}
\usepackage{tikz}
\usetikzlibrary{positioning,shapes.multipart,arrows.meta,shadows}
\usepackage[french]{babel}
\usepackage{xcolor}
\usepackage{listings}
\usepackage{graphicx}
\usepackage[most]{tcolorbox}

\definecolor{codegreen}{rgb}{0,0.6,0}
\definecolor{codegray}{rgb}{0.5,0.5,0.5}
\definecolor{codepurple}{rgb}{0.58,0,0.82}
\definecolor{backcolour}{rgb}{0.95,0.95,0.92}
 
\lstdefinestyle{mystyle}{
    backgroundcolor=\color{backcolour},   
    commentstyle=\color{codegreen},
    keywordstyle=\color{magenta},
    numberstyle=\tiny\color{codegray},
    stringstyle=\color{codepurple},
    basicstyle=\footnotesize,
    breakatwhitespace=false,         
    breaklines=true,                 
    captionpos=b,                    
    keepspaces=true,                 
    numbers=left,                    
    numbersep=5pt,                  
    showspaces=false,                
    showstringspaces=false,
    showtabs=false,                  
    tabsize=2
}
\lstset{style=mystyle, language=Python}

\title{CC1 1er NSI - Binaire, Variables et Conditions}
\author{Babash - NSI Tani Malandi et Johanne Bakalara}
\date{Octobre 2025, Durée: 2h}

\usepackage[a4paper,top=2.3cm,bottom=2.8cm,left=2cm,right=2cm,headsep=0.6cm,footskip=1.2cm]{geometry}
\lstset{extendedchars=true,literate={à}{{\`a}}1 {è}{{\`e}}1 {é}{{\'e}}1 {ê}{{\^e}}1}
\begin{document}
\runningfooter{}{\footnotesize Babash - NSI Tani Malandi et Johanne Bakalara \textperiodcentered{} CC BY-SA 4.0 \textperiodcentered{} page \thepage}{}
\firstpagefooter{}{\footnotesize Babash - NSI Tani Malandi et Johanne Bakalara \textperiodcentered{} CC BY-SA 4.0 \textperiodcentered{} page \thepage}{}

\maketitle
\thispagestyle{headandfoot}

\pagestyle{headandfoot}
\thispagestyle{headandfoot}
\firstpageheader{\makebox[\textwidth]{Nom : \dotfill\ Pr\'enom : \dotfill\ Classe : \makebox[3cm]{\dotfill}}}{}{}


\section*{Consignes:}
\begin{center} 
\fbox{\parbox{\linewidth}{
\begin{center}
Le devoir est à rédiger sur une feuille libre.\\
La clarté de la rédaction et les explications font partie intégrante de la notation. Une réponse non rédigée ou sans explication sera considérée fausse.\\
Prenez le temps d'encadrer les réponses.
\end{center}


}}
\end{center}


\vsword{Note}
\hsword{Note}

\section*{Barème :}
\begin{center}
  \par\bigskip
  \gradetable[h]
\end{center}


\section{Conversion Octet-bits: }

 Sam dispose d'un téléphone avec \textbf{80 Go} de mémoire. Chaque fichier musical pèse environ \textbf{240 Mbits}.
\begin{questions}
\question[2]
Combien de musiques complètes Sam peut-il stocker sur son téléphone ? (Vous détaillerez votre raisonnement et les conversions)

\section{Décimal, binaire et hexadécimal}

\begin{tcolorbox}[enhanced,
    colback=yellow!25!black!10!white,colframe=yellow!75!black,title=Rappel,
    drop fuzzy shadow,watermark color=white]
    \begin{itemize}
        \item On appelle la base 2 la base binaire, elle est composée uniquement de 2 symboles 0 et 1.
        \item On appelle la base 10 la base décimale, elle est composée de 10 symboles qui sont les chiffres de 0 à 9.
        \item On appelle la base 16 la base hexadécimale, elle est composée de 16 symboles qui sont les chiffres de 0 à 9 puis les lettres de A à F.
    \end{itemize}
\end{tcolorbox}

\question[3] Conversion binaire$\leftrightarrow$decimal
\begin{parts}
    \part \textbf{Convertissez} le nombre décimal $31_{10}$ \textbf{en binaire} (détaillez vos calculs).
    \part Pourquoi cette conversion était-elle 'facile' à anticiper ?
    \part \textbf{Convertissez} le nombre hexadécimal $A9_{16}$ \textbf{en décimal}.
\end{parts}

\begin{tcolorbox}[enhanced,
    colback=yellow!25!black!10!white,colframe=yellow!75!black,title=Rappel,
    drop fuzzy shadow,watermark color=white]
    % Table ASCII (codes 32 à 126), dessinée pour ce document.
\begin{center}
\footnotesize
\setlength{\tabcolsep}{3pt}
\renewcommand{\arraystretch}{1.3}
\begin{tabular}{|c|*{16}{c|}}
\hline
\textbf{hex} & \textbf{0} & \textbf{1} & \textbf{2} & \textbf{3} & \textbf{4} & \textbf{5} & \textbf{6} & \textbf{7} & \textbf{8} & \textbf{9} & \textbf{A} & \textbf{B} & \textbf{C} & \textbf{D} & \textbf{E} & \textbf{F} \\ \hline
\textbf{2} & \textsc{esp} & ! & " & \# & \$ & \% & \& & ' & ( & ) & * & + & , & - & . & / \\ \hline
\textbf{3} & 0 & 1 & 2 & 3 & 4 & 5 & 6 & 7 & 8 & 9 & : & ; & < & = & > & ? \\ \hline
\textbf{4} & @ & A & B & C & D & E & F & G & H & I & J & K & L & M & N & O \\ \hline
\textbf{5} & P & Q & R & S & T & U & V & W & X & Y & Z & [ & \textbackslash & ] & \^{} & \_ \\ \hline
\textbf{6} & ` & a & b & c & d & e & f & g & h & i & j & k & l & m & n & o \\ \hline
\textbf{7} & p & q & r & s & t & u & v & w & x & y & z & \{ & | & \} & \textasciitilde & \textsc{del} \\ \hline
\end{tabular}

\smallskip
Lecture : le caractère \texttt{A} est à la ligne \textbf{4}, colonne \textbf{1} : son code ASCII est \texttt{41} en hexadécimal ($65$ en décimal).
\end{center}

\end{tcolorbox}
\question[2] Chaque octet binaire ci-dessous représente un caractère ASCII. \textbf{Décodez le message} en le convertissant caractère par caractère (conseil : passez par l'hexadécimal).\\ \textbf{Message} :\\
01000010 01101001 01100101 01101110 00100000 01101010 01101111 01110101 01100101 00100000 00100001

\section{Complément à 2}

\question[1] Question de cours et réflexion
\begin{parts}
\part Indiquez si chacune des affirmations suivantes est VRAIE ou FAUSSE (sans justification demandée) :
\begin{itemize}
    \item En complément à 2 les nombres positifs gardent leurs représentations classiques.
    \item En complément à 2 l'addition d'un nombre et de son opposé vaut 0.
    \item En complément à 2 le zéro a deux représentations, une positive, une négative.
    \item En complément à 2 l'algorithme d'addition reste le même.
\end{itemize}
\end{parts}

\question[1] Sachant que le nombre binaire $1001 0111_{(2)}$ est codé en Complément à 2 sur 8 bits, déterminez sa valeur décimale.

\section{Variable}
\begin{lstlisting}[language=Python]
a = 6
b = 3
c = (a == 3 and b != 3)
b = a - b
b = a - b
e = a - (b+2)
a = e // 2
\end{lstlisting}
\question[1 \half] Considérez le bloc de code ci-dessus :
\begin{itemize}
    \item Combien de variables distinctes sont utilisées ?
    \item Quels sont leurs identifiants (noms) ?
    \item Quels sont les types de données finaux (int, float, etc.) de ces variables ?
\end{itemize}
\question[2] \textbf{Ecrivez le tableau de variables} pour l'exécution complète du programme, en indiquant la valeur de chaque variable après chaque ligne d'affectation.

\newpage
\section{Condition}

\begin{lstlisting}[language=Python]
# randint(1,6) tire au hasard une valeur entière comprise entre 1 et 6
de1 = randint(1,6)
de2 = randint(1,6)
de3 = randint(1,6)
if de1+1 == de2 and de2+1 == de3:
    score = 50
else :
    score = 0

if de1 == de2 == de3 :
    score = score + 200    
elif de1 == de2 :
    score = score + 20
elif de1 == de3 :
    score = score + 20
elif de2 == de3 :
    score = score + 20

if de1 == 6 or de2 == 6 or de3 == 6:
    score = score + 10

print("Vous marquez ",score," points")
\end{lstlisting}

\question[2] \textbf{Écrivez le tableau de variables} pour l'exécution complète du programme avec les valeurs initiales $de1=6$, $de2=6$ et $de3=6$.


\question[2] Indiquez uniquement la valeur finale de la variable \textit{score} pour les jeux de valeurs suivants (\textbf{aucune justification n'est requise}) :
\begin{itemize}
    \item $de1 = 6$, $de2 = 3$ et $de3 = 4$ :
    \item $de1 = 3$, $de2 = 4$ et $de3 = 5$ :
    \item $de1 = 3$, $de2 = 6$ et $de3 = 4$ :
    \item $de1 = 1$, $de2 = 1$ et $de3 = 4$ :
\end{itemize}

\newpage

\question[1 \half] Le programme doit vérifier si un mot de passe respecte trois critères simples. \textbf{Complétez le code} (...) pour vérifier ces critères.
\begin{lstlisting}[language=Python]
mot_de_passe_valide = True
longueur_mdp = 8      # Longueur du mot de passe
contient_chiffre = True # Le mot de passe contient-il un chiffre ?
contient_majuscule = False # Le mot de passe contient-il une majuscule ?

# Critère 1 : La longueur doit être supérieure ou égale à 12
if ...................................... :
    print("Critère 1 non respecté : Le mot de passe est trop court.")
    mot_de_passe_valide = False

# Critère 2 : Le mot de passe DOIT contenir un chiffre
if contient_chiffre == False :
    print("Critère 2 non respecté : Le mot de passe doit contenir un chiffre.")
    ...........................................

# Critère 3 : Le mot de passe DOIT contenir un chiffre ou une majuscule
if contient_majuscule == False and contient_chiffre == False :
    print("Critère 3 non respecté : Un mot de passe doit avoir soit un chiffre, soit une majuscule (ou les deux).")
    mot_de_passe_valide = False

if ...................................... :
    print("Mot de passe Validé.")
else:
    print("Mot de passe Non conforme.")
\end{lstlisting}

\question[2] \textbf{Écrivez un code Python} en utilisant les structures if, elif et else pour évaluer les variables \textit{temperature} et \textit{humidite} (saisies par l'utilisateur) et afficher le message approprié selon les règles suivantes :
\begin{itemize}
    \item "Journée idéale !" : Si la température est strictement entre 20 et 25 degrés (non inclus) et l'humidité est strictement inférieure à 15.
    \item "Chaleur lourde." : Si la température est supérieure ou égale à 25 et l'humidité est strictement supérieure à 15.
    \item "Prévoyez une veste." : Si la température est inférieure ou égale à 15 degrés.
    \item "Météo agréable." : Dans tous les autres cas.
\end{itemize}
La \textit{temperature} et l'\textit{humidite} seront saisies par l'utilisateur, vous pourrait utiliser :
\begin{lstlisting}[language=Python]
temperature = int(input("Saisissez la température :"))
humidite = int(input("Saisissez l'humidité :"))
\end{lstlisting}

\end{questions}
\end{document}

